% =============================================================================
% Kapitel 6: Diskussion
% -----------------------------------------------------------------------------
% Was hinein gehört (Richtwert 5--8 Seiten):
%  - Beantwortung der Forschungsfragen F1--F4 aus Kapitel 1, je ein Abschnitt,
%    mit Verweis auf die konkreten Ergebnisse (\cref{tab:...}, \cref{fig:...})
%  - Einordnung gegenüber dem Stand der Technik (Kapitel 3): Was ist besser,
%    was schlechter, warum?
%  - Genauigkeit vs. Laufzeit: Wann lohnt sich das Netz? Hybrid (Netz als
%    Startwert für die modellbasierte Optimierung)?
%  - Grenzen: Simulationsannahmen, Sim-to-Real-Gap, Verteilungsverschiebung,
%    Abhängigkeit von z02/Energie/Pixelgröße (Generalisierung über Geometrien),
%    Datensatzgröße, Rechenaufwand
%  - Bedrohungen der Validität: Datenleckage zwischen Splits, Seeds, Auswahl der
%    Metriken, Referenz-Fr der Messdaten selbst nur geschätzt
% =============================================================================
\chapter{\lang{Diskussion}{Discussion}}
\label{ch:diskussion}

\section{\lang{Beantwortung der Forschungsfragen}{Answering the Research Questions}}
\label{sec:diskussion:fragen}
\todo{\lang{F1--F4 mit Verweis auf \cref{tab:experimente:baseline} usw.}{F1--F4 with reference to \cref{tab:experimente:baseline} etc.}}

\section{\lang{Einordnung gegenüber dem Stand der Technik}{Comparison with the State of the Art}}
\label{sec:diskussion:einordnung}
\todo{\lang{Vergleich mit \cite{dora2025autofocus} und den Arbeiten aus \cref{ch:stand}.}{Comparison with \cite{dora2025autofocus} and the works from \cref{ch:stand}.}}

\section{\lang{Grenzen der Arbeit}{Limitations}}
\label{sec:diskussion:grenzen}
\todo{\lang{Simulationsannahmen, Sim-to-Real-Gap, Generalisierung über Geometrien.}{Simulation assumptions, sim-to-real gap, generalisation across geometries.}}
